\subsection{Uniform equations cut out by marked torsion}
\label{subsec:ftb-uniform-torsion-forms}

For a finite set \(\Sigma\), let
\(A_\Sigma(K)=K[X_i:i\in\Sigma]\) and
\[
                         m_\Sigma(e)=\dim_K A_\Sigma(K)_e .
\]
For nonempty \(\Sigma\) this is the value at \(e\ge0\) of
one polynomial \(M_\Sigma\in\mathbb Q[T]\), independently
of \(K\).
We use the source's ideal-valued Hilbert function and the
following literal condition.
For vectors \(p_h\in K^\Sigma\), define
\[
 V_d(p)=\bigl(F:F\in A_\Sigma(K)_d,\
                                  F(p_h)=0\text{ for all }h\bigr).
\]
The assertion \(\operatorname{CutBy}(I,Q,d,p)\) means
\[
       \dim_K I_e=Q(e),\qquad
       \bigl(F\in I\Longleftrightarrow F\in V_d(p)\bigr)
                  \quad(F\in A_\Sigma(K)_e,\ e\ge d).
 \tag{TF1}\label{eq:ftb-cutby-definition}
\]
Thus \(V_d(p)\) is the ideal generated by the degree-\(d\)
forms, rather than merely their vector space.

We first verify the actual ideal to which regularity is applied.
Let \(Y/K\) have a line bundle \(M\) and a finite very ample
system \((\theta_i)\).
A section-coordinate vector at \(y:\operatorname{Spec}R\to Y\)
is obtained by expressing the pulled-back sections in one
generator of \(y^*M\).
The ideal \(I(\theta)\) consists of polynomials vanishing on
every such vector over every \(K\)-algebra \(R\).
This retains the scheme-valued, including nonreduced,
test-ring definition of the source.

\begin{lemma}[The actual embedding ideal]
\label{lem:ftb-theta-embedding-ideal}
The ideal \(I(\theta)\) is homogeneous and satisfies
\[
                   (X_iF\in I(\theta)\text{ for all }i)
                                      \Longrightarrow F\in I(\theta).
 \tag{TF2}\label{eq:ftb-embedding-ideal-saturated}
\]
A homogeneous form \(F\) of degree \(e\) defines the actual
section \(F(\theta)\in H^0(Y,M^e)\).
It is zero if and only if \(F\in I(\theta)\), and at a
section-coordinate vector \(v\) over a field it vanishes
if and only if \(F(v)=0\).
For all sufficiently large \(e\),
\[
            \dim_K I(\theta)_e+h^0(Y,M^e)=m_\Sigma(e).
 \tag{TF3}\label{eq:ftb-actual-ideal-hilbert}
\]
\end{lemma}
\begin{proof}
Scale an arbitrary section-coordinate vector \(v\) by
the invertible variable \(Z\) over \(R[Z,Z^{-1}]\).
If \(F\) vanishes on every such vector, then
\(F(Zv)=0\).
Its coefficient of \(Z^e\) is the value of the
degree-\(e\) homogeneous component of \(F\) at \(v\).
This proves homogeneity over all test rings.

The coordinates of any section-coordinate vector
generate the unit ideal.
Otherwise a prime containing all of them would give a
point of \(Y\) where every \(\theta_i\) vanished, contrary
to the covering property of the system.
A finite linear combination of its coordinates equal
to \(1\) therefore cancels the common factor \(F(v)\)
from the equations \(v_iF(v)=0\).
This proves \eqref{eq:ftb-embedding-ideal-saturated}
without passing to reduced test rings.

In a trivializing chart of \(M\), substitute the actual
functions of the \(\theta_i\) in \(F\).
Homogeneity makes their changes of trivialization the
\(e\)-th powers of those of \(M\), so these functions
glue to \(F(\theta)\).
At any test point, its local function is the value
\(F(v)\) multiplied by the \(e\)-th power of the
chosen invertible generator.
This proves the two vanishing equivalences.
For the reverse implication from all test vectors to
the zero section, use affine subcharts of a
trivializing chart and their identity-valued points;
the same unit factor can be cancelled there.

The very ample system is the specified projective
coordinate embedding with coordinate line \(M\).
Serre vanishing for its coherent ideal sheaf gives,
for all sufficiently large \(e\), the surjection
\[
                        A_\Sigma(K)_e\longrightarrow H^0(Y,M^e).
\]
Its kernel is exactly \(I(\theta)_e\) by the previous
paragraph.
The vector-space kernel/image dimension identity gives
\eqref{eq:ftb-actual-ideal-hilbert}.
This is the source's eventual section-surjectivity step;
no uniform bound for that surjectivity threshold is
asserted or needed.
\end{proof}

\begin{proposition}[The source's uniform torsion-form bound]
\label{prop:ftb-uniform-torsion-forms}
Fix \(g,C\in\mathbb N\) and a nonempty finite coordinate
set \(\Sigma_0\).
There are integers \(d\ge1\) and \(k_0\), before all
fields, abelian varieties, and repeated coordinate
systems below, with the following property.
Let \(r:\Sigma\twoheadrightarrow\Sigma_0\) be a
surjection of finite sets.
There is a polynomial \(Q\), depending only on
\(g,C,\Sigma_0,\Sigma,r,\delta\), for each
\(0\le\delta\le C^g\), such that the following data
satisfy \eqref{eq:ftb-cutby-definition}:
\begin{enumerate}
\item an algebraically closed field \(K\), an abelian
variety \(A/K\) of dimension \(g\), an isomorphism
\(Y\simeq A\), and a line bundle \(M\) on \(Y\);
\item the transported line bundle on \(A\) is symmetric,
and \((\theta_i)_{i\in\Sigma}\) is a very ample system
for \(M\), with
\(\theta_i=\theta_{i'}\) whenever \(r(i)=r(i')\);
\item the exact section dimensions are
\[
                       h^0(Y,M^e)=\delta e^g\quad(e\ge1);
 \tag{TF4}\label{eq:ftb-exact-section-growth}
\]
\item \(k\ge k_0\), \(k\ne0\) in \(K\), a family
\(y_h\in A(K)\) containing every point of \(A[k](K)\),
and actual section-coordinate vectors \(p_h\) of
\((\theta_i)\) at their corresponding points of \(Y\).
\end{enumerate}
The conclusion is
\[
                         \operatorname{CutBy}(I(\theta),Q,d,p).
 \tag{TF5}\label{eq:ftb-uniform-torsion-cut}
\]
One may take precisely
\[
 \begin{split}
 d&=1+\max_{0\le\delta\le C^g}
                   d_0\bigl(\Sigma_0,M_{\Sigma_0}-\delta T^g\bigr),\\
 k_0&=C^g g!d+1,\qquad
                         Q(T)=M_\Sigma(T)-\delta T^g,
 \end{split}
 \tag{TF6}\label{eq:ftb-torsion-cut-choices}
\]
where \(d_0\) is from Lemma~\ref{lem:ftb-uniform-regularity}.
The conclusion also holds when \(g=0\).
\end{proposition}
\begin{proof}
Choose a right inverse \(s\) to \(r\).
The repeated sections descend to a system
\(\theta^0_j=\theta_{s(j)}\), with
\(\theta_i=\theta^0_{r(i)}\).
The smaller system is still very ample: the same
nonzero-locus charts cover \(Y\); their ratios generate
the affine coordinate rings, since repeated ratios add
no functions.
Writing \(I_0=I(\theta^0)\) and \(I=I(\theta)\), the
all-test-ring definition gives
\[
                  I=(\operatorname{rename}r)^{-1}I_0.
 \tag{TF7}\label{eq:ftb-repeated-coordinate-ideal}
\]
Indeed, section-coordinate vectors of the repeated
system are repeated vectors.
Conversely their actual common line generator also
gives the vector for the smaller system, and equality
of the repeated coordinates follows by cancellation
of that generator.

Lemma~\ref{lem:ftb-theta-embedding-ideal} and
\eqref{eq:ftb-exact-section-growth} give the eventual
ideal Hilbert polynomial
\(M_{\Sigma_0}-\delta T^g\).
Its ideal is homogeneous and saturated in the exact
sense required by uniform regularity.
Therefore, for every \(e\ge d\), its degree-\(e\)
part has the prescribed dimension and is generated by
its degree-\(d\) part.
To see why the source can raise the initially obtained
degree \(d_0\) to \(d\), expand homogeneous polynomial
multipliers into monomials.
In a product of total degree \(e\ge d\), split off
\(d-d_0\) variables of the multiplier.
Those variables times the original degree-\(d_0\)
generator form a degree-\(d\) element of the ideal.
Thus generation persists at every larger chosen degree.

Here is the exact repeated-coordinate transfer, so no
bound depends on \(\#\Sigma\) through a new choice of
\(d\).
The maps \(\operatorname{rename}r\) and
\(\operatorname{rename}s\) split on every homogeneous
piece.
Their kernels on the full piece and on \(I_e\) agree.
Subtracting their two kernel/image dimension formulas
gives
\[
                    \dim_K I_e
                  =\dim_K(I_0)_e+m_\Sigma(e)-m_{\Sigma_0}(e).
 \tag{TF8}\label{eq:ftb-repeated-coordinate-dimension}
\]
For generation let \(t=s\circ r\).
The kernel differences \(X_i-X_{t(i)}\) lie in \(I_1\).
For \(e\ge d\ge1\), the difference
\(F-\operatorname{rename}t(F)\), with \(F\) homogeneous
of degree \(e\), belongs to the ideal generated by \(I_d\).
This follows by induction on a product of linear
forms, using
\[
 lG-(\operatorname{rename}t(l))
              (\operatorname{rename}t(G))
 =l(G-\operatorname{rename}t(G))
  +(l-\operatorname{rename}t(l))\operatorname{rename}t(G).
\]
In the second term, split \(d-1\) factors from the
degree-\((e-1)\) multiplier; the first term uses
induction.
The remaining term
\(\operatorname{rename}t(F)\) is generated by \(I_d\)
by pulling back the already proved generation for \(I_0\).
This proves both the dimension polynomial \(Q\) in
\eqref{eq:ftb-torsion-cut-choices} and generation for \(I\).

It remains to identify the degree-\(d\) equations from
the marked points.
Every \(F\in I\) vanishes at every \(p_h\), by its
definition.
Conversely suppose \(F\in A_\Sigma(K)_d\) vanishes at
every \(p_h\).
Its actual section \(\tau=F(\theta)\) vanishes at all
of \(A[k](K)\).
If \(g\ge1\) and \(\tau\ne0\), transport it and \(M\)
along \(Y\simeq A\) and apply
Proposition~\ref{prop:ftb-torsion-zero-growth}.
The dimension assumption gives
\(h^0(A,M^e)=\delta e^g\le(Ce)^g\), so
\[
                       k^{2g}\le C^g g!d\,k^{2g-2}.
\]
As \(k\ge1\), this implies \(k^2\le C^g g!d\),
contrary to \(k>C^g g!d\).
For \(g=0\), the origin is dense in the integral
zero-dimensional abelian variety; vanishing there
already makes \(\tau=0\).
In all cases the embedding-ideal lemma gives \(F\in I\).
Thus \(V_d(p)\) is exactly the ideal generated by \(I_d\).
The generation and dimension conclusions above prove
\eqref{eq:ftb-uniform-torsion-cut}.
\end{proof}

\paragraph{Source concordance and next application.}
The scheme-valued ideal and polynomial-section proof are
\texttt{ModuliProofs/HomogeneousIdealSectionsOfLineBundlePower.lean}
(\texttt{e96c0f4620fbc0b399798b6bec2a0f1e95e6f0c3});
the actual eventual-surjectivity and kernel/image calculation are
\texttt{HilbertFunctionVanishingIdealVeryAmpleSystem.lean}
(\texttt{a0e221649a802760557fdfefb84ac5a603b86fdb}).
The repeated-coordinate transfer is
\texttt{FinrankMapAddFinrankInfKerOfFinite.lean}
(\texttt{a426bd2cbe4777259301e1cbc787af43926696ab});
the degree-raising and final \(\operatorname{CutBy}\) assembly are
\texttt{HomogeneousComponentMulMemSpanDegree.lean}
(\texttt{50a8605d81d8024127d5168dadc5d65f40e82eb7}).
The literal final producer, including the choices
\eqref{eq:ftb-torsion-cut-choices}, is
\texttt{AbelianVarietyCutOutBoundedDegreeTorsion.lean}
(\texttt{99d60175c39a88b37637675c504a372e3f2aa332}).
All paths have prefix
\texttt{src/Development/AbelianVarietyTheory/Proofs/}.
This closes its full supplied-system scope.
To obtain the actual fine Siegel family one still has
to instantiate the very ample, periodic theta systems
and the addition/inversion/endomorphism graph systems.
The preceding proposition does not supply those systems
by itself.
